\chapter*{Summary}
\addcontentsline{toc}{chapter}{Summary}


The computational cost of scale-resolving CFD is dominated, in most mesh-based incompressible solvers, by the time integration of the discretised Navier--Stokes equations and the pressure projection that enforces incompressibility at every step. This thesis investigates how far a learned latent-space surrogate can replace that step in an unsteady, vortex-shedding flow without degrading the physical fidelity of the simulation it accelerates.


The proposed framework couples a convolutional autoencoder, a latent neural ordinary differential equation (NODE), and a runtime out-of-distribution (OOD) gate to the differentiable incompressible solver WaterLily.jl. The autoencoder compresses each instantaneous velocity field, augmented with the boundary-data-immersion geometry field, into a 16-dimensional latent state, and is trained with a physics-informed objective that adds divergence and vorticity penalties to the $L_1$ reconstruction loss, so the decoded fields stay near-divergence-free and retain the shear-layer vorticity that sets the near-body loading. The NODE learns the latent dynamics and is fitted with multiple shooting, partitioning each trajectory into short segments to suppress the gradient blow-up that destabilises single-shooting training over the long, chaotic rollouts spanning several shedding cycles. At inference, a $k$-nearest-neighbour (KNN) monitor measures the distance between each predicted latent vector and the encoded training set; once a rollout drifts beyond a calibrated threshold, the prediction is rejected, and control returns to WaterLily, which re-stabilises the flow before the loop resumes. Repeated flags trigger retraining of the NODE on the newly gathered trajectories.


The framework is developed on two-dimensional flow past a circular cylinder at $Re = 2500$ and evaluated on $256^2$ and $512^2$ grids. In its in-distribution regime, the hybrid loop preserves both the engineering- and statistics-level observables: on the finer grid, the time-averaged drag is recovered to within $0.64\%$ over the accelerated windows, and the mean-velocity and Reynolds-stress fields correlate spatially with the reference at $\rho \geq 0.997$, with the fluctuating lift the only quantity to carry a visible rollout error (its r.m.s.\ amplitude decaying to a roughly $7\%$ error inside the hybrid regions). Once trained, a single latent rollout is about $85\times$ cheaper per convective time unit (CTU) than the reference time integration, giving an inference-loop speedup of roughly $1.48\times$ at the base resolution and climbing above $2\times$ on the refined $512^2$ grid, since the memory-bound solver cost scales with grid size while the rollout cost is fixed by the latent dimension.


The decisive limitation is economic rather than physical. Charging the one-off network training to the comparison collapses the effective speedup of a from-scratch run to roughly $0.07\times$ at the base point, about an order of magnitude slower than simply running WaterLily, because training consumes the large majority of the wall-clock budget and is not repaid within a single simulation window. The parameter studies expose one recurring trade-off behind this: narrowing the latent bottleneck, shortening the autoencoder schedule, or loosening the multiple-shooting group each simplifies the dynamics the NODE must learn and lengthens the average in-distribution rollout, but coarsens the recovered second-order statistics, so acceleration and turbulence-statistics fidelity are pulled against one another, with the OOD gate setting the operating point on that curve. The training cost is not irreducible: it amortises across reused pipelines, and trimming the schedule alone lifts the effective speedup to about $0.62$ on the finer grid. Pushed out of distribution to $Re = 5000$, the autoencoder smooths the held-out near-wall fluctuations and the instantaneous loading becomes unreliable (hybrid-region force errors approaching $9\%$) even as the time-averaged fields stay faithful, locating the autoencoder's reconstruction ceiling as the binding accuracy constraint. The work thus establishes that latent NODE rollouts can absorb part of the conventional time integration without sacrificing the integrated loading or the wake statistics, and identifies training cost, not the inference loop, as the bottleneck any practical deployment must overcome.
